% =============================================================
% Section 2: Background and Related Work
% Target: ~1.5 columns (~750 words in 2-col AAAI format)
% =============================================================

\section{Background and Related Work}
\label{sec:background}

\subsection{Time Series Foundation Models}
\label{sec:bg_tsfm}

The past two years have produced a distinct class of models trained on large,
heterogeneous collections of time series with the aim of zero-shot or few-shot
transfer across domains.
These models vary along two architectural axes: tokenisation strategy and
sequence-to-sequence objective.

\textbf{Patch-based encoder-decoder models.}
MOMENT~\cite{goswami2024moment} applies masked patch modelling to sub-series
segments, learning a representation that can be adapted to forecasting,
classification, and anomaly detection.
Moirai~\cite{woo2024moirai} extends this with a multi-patch-size architecture
that trains a single model to handle variable frequencies.

\textbf{Decoder-only models.}
TimesFM~\cite{das2024timesfm} treats forecasting as next-patch prediction,
training a 500M-parameter decoder-only model on a corpus spanning Google
Trends, Wikipedia pageviews, and synthetic series.
Timer~\cite{liu2024timer} follows a similar autoregressive design at smaller scale.

\textbf{Language-model-inspired models.}
Chronos~\cite{ansari2024chronos} quantises time series values into discrete tokens
and fine-tunes T5~\cite{raffel2020t5} encoder-decoder architectures on a curated
public corpus, treating forecasting as a sequence-to-sequence language task.
Lag-Llama~\cite{rasul2024lagllama} adapts LLaMA with lag features as covariates.

Comprehensive surveys of this space are provided by
\citet{jin2023large} and \citet{liang2024foundation}.
Despite broad empirical coverage, neither survey addresses the internal geometric
structure of learned representations---the question this paper targets.

\textbf{The case against complexity.}
Before attributing zero-shot transfer to any internal property of TSFMs,
a direct challenge must be addressed.
\citet{zeng2023transformers} demonstrated that a simple one-layer linear model
matches or outperforms Transformer-based forecasters on five standard long-horizon
benchmarks, including ETT and Weather.
This result does not directly refute TSFMs---linear models were not tested in
zero-shot settings, and the comparison is against models trained per-dataset---but
it establishes that Transformer complexity is neither necessary nor sufficient for
forecasting quality in supervised settings.
Our geometric analysis is designed with this sceptical position in mind:
rather than asserting TSFMs are better, we ask whether their latent representations
are geometrically distinct from supervised baselines trained on the same data,
and whether any such distinction correlates with forecast reliability.

\subsection{Geometric Analysis of Neural Representations}
\label{sec:bg_geometry}

The internal geometry of neural network representations has been studied
systematically in vision and language models, establishing a set of
analytical tools directly applicable to our setting.

\textbf{Intrinsic dimensionality.}
\citet{ansuini2019intrinsic} showed that deep image classifiers exhibit a
characteristic ID profile: representations expand to high effective
dimensionality in early layers and compress toward a low-dimensional manifold
in the final layers, with compression correlating with generalisation quality.
\citet{birdal2021intrinsic} extended this to language models, finding that
low intrinsic dimensionality of the fine-tuning landscape explains why
large pretrained models are efficiently adapted with few parameters.
We apply the same analysis to TSFMs, where the relevant question is whether
the compression pattern reflects meaningful temporal abstraction or merely
the inductive bias of the Transformer architecture.

\textbf{Geometry of language model representations.}
\citet{reif2019visualizing} applied geometric analysis to BERT, finding that
contextual embeddings form geometrically separable clusters corresponding to
syntactic and semantic categories.
\citet{kornblith2019similarity} introduced Centred Kernel Alignment (CKA) to
compare representation geometry across architectures, showing that wider and
deeper networks converge to similar representations in intermediate layers.
No equivalent geometric analysis exists for time series models.

\textbf{Representational similarity.}
\citet{raghu2017svcca} introduced SVCCA to track representation geometry through
training, demonstrating that lower layers stabilise early while upper layers
continue evolving---a pattern relevant to our layer-wise ID profile analysis.

\subsection{Explainability for Time Series}
\label{sec:bg_xai_ts}

Existing XAI methods for time series fall into two categories:
attribution methods that assign importance scores to input timesteps
(LIME, SHAP, integrated gradients), and perturbation-based methods that
measure model sensitivity to input modifications~\cite{ismail2020benchmarking,
theissler2022explainable}.
Both categories operate at the input level and produce post-hoc explanations
for individual predictions.
Neither provides a structural account of what a model has learned to encode.

Our approach is fundamentally different: we analyse the \emph{representation space}
itself---its dimensionality, stability, and semantic organisation---rather than
attributing individual predictions.
This distinction aligns GEF with mechanistic analysis rather than attribution-based
XAI, and the two approaches are complementary rather than competing.

\subsection{Linear Probing and Temporal Primitives}
\label{sec:bg_probing}

Linear probing~\cite{alain2017probing} tests whether a target property is linearly
decodable from frozen intermediate representations, providing a controlled
measure of what information is encoded without confounds from end-to-end training.
In NLP, probing has revealed that BERT encodes syntactic structure in intermediate
layers and semantic structure in upper layers~\cite{reif2019visualizing}.
We apply this methodology to test whether TSFM embeddings encode three
interpretable temporal primitives: trend direction, volatility regime, and
seasonality---properties that are trivially computable from the input but
non-trivially encoded if the representation space organises them as separable
geometric subspaces.

\subsection{Theoretical Motivation}
\label{sec:bg_theory}

A classical result in dynamical systems theory offers an intuitive lens.
Takens' embedding theorem~\cite{takens1981detecting} states that delay embeddings
of scalar observations from a smooth deterministic dynamical system reconstruct
the system's attractor manifold, up to diffeomorphism.
In other words, finite windows of a time series generated by a dynamical system
already lie on a low-dimensional manifold in delay-embedding space.

We invoke this result as \emph{motivating intuition only}.
The theorem's conditions---smooth, noise-free, deterministic dynamics---are not
satisfied by the heterogeneous real-world series (electricity consumption, traffic
speed, meteorological observations) studied here.
Our hypothesis, stated without claiming the theorem's authority, is that TSFMs
trained on diverse corpora have learned to approximate the manifold structure of
temporal data in their latent space, and that this approximation is measurable
via the geometric tools described in Section~\ref{sec:framework}.
Whether this constitutes genuine manifold learning or a looser form of
distributional compression is precisely what our experiments test.
